\subproblem{Дыбыс интерференциясы}{3.0}

Екі шексіз қатты жазықтық (қабырғалар) тік бұрыш жасап қиылысады. Шексіздіктен жазық дыбыс толқыны келіп түседі, оның толқындық векторы дәл жазықтықтар арасындағы бұрыштың биссектрисасы бойымен бағытталған. Түскен толқындағы ауа молекулалары жылдамдығының амплитудасы $v$, ал толқын ұзындығы $\lambda$-ға тең. Қабырғалар идеал шағылыстырады (шағылысу энергия шығынынсыз жүреді); $\lambda$ ретіндегі қашықтықтағы дифракциялық эффектілерді ескермеңіз.

\setcounter{taskstar}{0}
\begin{list}{\textbf{(\alph{taskstar})}\ }{
\usecounter{taskstar}
\setlength{\leftmargin}{18pt}%
\setlength{\labelwidth}{0pt}%
\setlength{\labelsep}{0pt}%
\setlength{\itemsep}{6pt}%
\setlength{\topsep}{6pt}}
\item Осы жүйедегі ауа қозғалысының мүмкін болатын максималды жылдамдығын табыңыз.
\item Кеңістіктің қай нүктелерінде осы максималды жылдамдыққа қол жеткізіледі?
\end{list}
